\section{The trijet cross-section}
\label{sec:xs}

We now insert the outgoing state \eqref{eq:state} into eq.~\eqref{eq:xsdef}. The number-density
operators are expressed in the transverse coordinate representation,
\begin{equation}
N_q(k^+,\kk)=\frac{1}{(2\pi)^3}\int_{\x,\bx}e^{-i\kk\cdot(\x-\bx)}\;
b^{\alpha\dagger}_\lambda(k^+,\bx)\,b^{\alpha}_\lambda(k^+,\x)\,,
\end{equation}
and similarly for the antiquark and the gluon, and the Fock-space operators are contracted with those
of the state. The integrations over $\w$ and $\bar\w$ remove the $\delta$-functions which identify the
photon position with the centre of energy $\cc$. The cross-section is the sum of three terms: the
direct contributions of the antiquark and of the quark, in which the gluon is emitted by the same
fermion in the direct amplitude (DA) and in the complex-conjugate amplitude (CCA), and the
interference contribution,
\begin{equation}
\frac{\dd\sigma^{\gT A\to\qqg+X}}{\dd^3k_1\,\dd^3k_2\,\dd^3k_3}
=\frac{\dd\sigma^{(\bar q\bar q)}}{\dd^3k_1\,\dd^3k_2\,\dd^3k_3}
+\frac{\dd\sigma^{(qq)}}{\dd^3k_1\,\dd^3k_2\,\dd^3k_3}
+\frac{\dd\sigma^{(q\bar q)}}{\dd^3k_1\,\dd^3k_2\,\dd^3k_3}\,.
\label{eq:xs-sum}
\end{equation}
The longitudinal fractions are fixed by the final state, $\zeta=k_1^+/q^+$, $\xi=k_3^+/q^+$ and
$\eta=k_2^+/q^+=1-\zeta-\xi$, and are therefore the same in the DA and in the CCA. Coordinates and
functions with a bar refer to the CCA: for instance $\bR=\bx-\byp$, $\bY=\by-\bz$ and
$\bar F^{(\alpha\beta)}\equiv F^{(\alpha\beta)}(\bR,\bY)$. Since each of the amplitudes
\eqref{eq:state-qbar} and \eqref{eq:state-q} is the difference of two terms, each of the three
contributions in eq.~\eqref{eq:xs-sum} contains four terms.

\subsection{Colour structure}
\label{sec:colour}

The average over the colour fields of the target generates the $S$-matrices of colourless systems made
of up to six partons. For the terms in which the three partons scatter both in the DA and in the CCA we
need
\begin{align}
\cS_{\qqg\,\qqg}(\x,\y,\z;\bx,\by,\bz)&\equiv\frac{1}{\CF\Nc}\Big\langle\big[U^\dagger(\bz)U(\z)\big]^{ab}\;
\tr\big[V(\by)\,t^a\,V^\dagger(\bx)\,V(\x)\,t^b\,V^\dagger(\y)\big]\Big\rangle\,,
\nonumber\\
\cS_{\qqg}(\x,\y,\z)&\equiv\frac{1}{\CF\Nc}\Big\langle U^{ab}(\z)\;
\tr\big[V(\x)\,t^b\,V^\dagger(\y)\,t^a\big]\Big\rangle\,,
\label{eq:S6}
\end{align}
where the brackets denote the CGC average. The trace of the product of the Wilson-line structure of
eq.~\eqref{eq:wf-a} with its complex conjugate is then $\CF\Nc\,\cW$, with
\begin{equation}
\cW(\x,\y,\z;\bx,\by,\bz)\equiv\cS_{\qqg\,\qqg}(\x,\y,\z;\bx,\by,\bz)-\cS_{\qqg}(\x,\y,\z)
-\cS^\dagger_{\qqg}(\bx,\by,\bz)+1\,.
\label{eq:W}
\end{equation}
These expressions hold for any $\Nc$. Since the Wilson-line structures for an emission after the
collision are the coincidence limits of the one in eq.~\eqref{eq:wf-a}, namely $\y,\z\to\yp$ for the
antiquark and $\x,\z\to\xp$ for the quark, all the colour structures that we need are obtained from
eq.~\eqref{eq:W} by taking the corresponding limits, in the DA and/or in the CCA.

At large $\Nc$ the adjoint Wilson line is the product of a fundamental and an anti-fundamental one,
and eq.~\eqref{eq:S6} reduces to products of dipoles and quadrupoles,
\begin{align}
S(\x,\y)&\equiv\frac{1}{\Nc}\Big\langle\tr\big[V(\x)V^\dagger(\y)\big]\Big\rangle\,,
\nonumber\\
\cQ(\x_1,\x_2,\x_3,\x_4)&\equiv\frac{1}{\Nc}\Big\langle\tr\big[V(\x_1)V^\dagger(\x_2)V(\x_3)V^\dagger(\x_4)\big]\Big\rangle\,,
\end{align}
namely $\cS_{\qqg\,\qqg}\simeq\cQ(\x,\z,\bz,\bx)\,\cQ(\z,\y,\by,\bz)$ and
$\cS_{\qqg}\simeq S(\x,\z)\,S(\z,\y)$~\cite{Iancu:2022gpw}. Hence
\begin{equation}
\cW(\x,\y,\z;\bx,\by,\bz)\simeq\cQ(\x,\z,\bz,\bx)\,\cQ(\z,\y,\by,\bz)-S(\x,\z)\,S(\z,\y)
-S(\bz,\bx)\,S(\by,\bz)+1\,.
\label{eq:W-largeN}
\end{equation}
The coincidence limits which are relevant for the emission by the antiquark are
\begin{align}
\cW(\x,\yp,\yp;\bx,\by,\bz)&\simeq\cQ(\x,\yp,\bz,\bx)\,S(\by,\bz)-S(\x,\yp)-S(\bz,\bx)\,S(\by,\bz)+1\,,
\nonumber\\
\cW(\x,\y,\z;\bx,\byp,\byp)&\simeq\cQ(\x,\z,\byp,\bx)\,S(\z,\y)-S(\x,\z)\,S(\z,\y)-S(\byp,\bx)+1\,,
\nonumber\\
\cW(\x,\yp,\yp;\bx,\byp,\byp)&\simeq\cQ(\x,\yp,\byp,\bx)-S(\x,\yp)-S(\byp,\bx)+1\,,
\label{eq:W-qbar}
\end{align}
those for the emission by the quark are
\begin{align}
\cW(\xp,\y,\xp;\bx,\by,\bz)&\simeq S(\bz,\bx)\,\cQ(\xp,\y,\by,\bz)-S(\xp,\y)-S(\bz,\bx)\,S(\by,\bz)+1\,,
\nonumber\\
\cW(\x,\y,\z;\bxp,\by,\bxp)&\simeq S(\x,\z)\,\cQ(\z,\y,\by,\bxp)-S(\x,\z)\,S(\z,\y)-S(\by,\bxp)+1\,,
\nonumber\\
\cW(\xp,\y,\xp;\bxp,\by,\bxp)&\simeq\cQ(\xp,\y,\by,\bxp)-S(\xp,\y)-S(\by,\bxp)+1\,,
\label{eq:W-q}
\end{align}
and the interference terms involve in addition
\begin{align}
\cW(\xp,\y,\xp;\bx,\byp,\byp)&\simeq S(\byp,\bx)\,S(\xp,\y)-S(\xp,\y)-S(\byp,\bx)+1
\nonumber\\
&=\big[S(\xp,\y)-1\big]\big[S(\byp,\bx)-1\big]\,.
\label{eq:W-int}
\end{align}
The last lines of eqs.~\eqref{eq:W-qbar} and \eqref{eq:W-q} have the same structure as the colour
factor of the leading-order dijet cross-section,
\begin{equation}
\cW_{\rm LO}(\x,\y;\bx,\by)\equiv\cQ(\x,\y,\by,\bx)-S(\x,\y)-S(\by,\bx)+1\,,
\label{eq:WLO}
\end{equation}
as expected for the graphs in which the gluon is emitted in the final state on both sides of the cut.

The dipole $S$-matrix can be taken from a solution of the BK equation or from a model. The quadrupole
is in general an independent correlator, but in the Gaussian approximation and at large
$\Nc$~\cite{Dominguez:2011wm,Iancu:2011ns,Iancu:2011nj} it can be expressed in terms of the dipole,
\begin{align}
\cQ(\x_1,\x_2,\x_3,\x_4)\simeq{}& S(\x_1,\x_2)\,S(\x_3,\x_4)
\nonumber\\
&-\frac{F(\x_1,\x_2;\x_3,\x_4)}{F(\x_1,\x_3;\x_2,\x_4)}\,
\Big[S(\x_1,\x_2)\,S(\x_3,\x_4)-S(\x_1,\x_4)\,S(\x_3,\x_2)\Big]\,,
\label{eq:Gauss}
\end{align}
with
\begin{equation}
F(\x_1,\x_2;\x_3,\x_4)\equiv\ln\frac{S(\x_1,\x_3)\,S(\x_2,\x_4)}{S(\x_1,\x_4)\,S(\x_2,\x_3)}\,.
\label{eq:F}
\end{equation}
Equation~\eqref{eq:Gauss} has the correct coincidence limits, $\cQ(\x_1,\x_2,\x_2,\x_4)=S(\x_1,\x_4)$ and
$\cQ(\x_1,\x_2,\x_3,\x_1)=S(\x_3,\x_2)$. With eqs.~\eqref{eq:W-largeN}--\eqref{eq:Gauss} the whole trijet
cross-section is expressed in terms of the dipole $S$-matrix alone.

\subsection{Direct contribution of the antiquark}
\label{sec:xs-qbar}

Using eq.~\eqref{eq:state-qbar} and $e^2g^2=(4\pi)^2\aem\as$, we find
\begin{align}
\frac{\dd\sigma^{(\bar q\bar q)}}{\dd^3k_1\,\dd^3k_2\,\dd^3k_3}=&\;
\frac{\aem\,\as\,\CF\Nc}{2(2\pi)^{10}}\;\frac{\zeta^2}{\xi\,(q^+)^2}\;
\delta\big(q^+-k_1^+-k_2^+-k_3^+\big)\sum_fe_f^2
\nonumber\\
&\times\int_{\x,\y,\z,\bx,\by,\bz}e^{-i\kk_1\cdot(\x-\bx)-i\kk_2\cdot(\y-\by)-i\kk_3\cdot(\z-\bz)}\;
\Big[\cH[\cG;\bar\cG]\;\cW(\x,\y,\z;\bx,\by,\bz)
\nonumber\\
&\qquad-\cH[\cA;\bar\cG]\;\cW(\x,\yp,\yp;\bx,\by,\bz)
-\cH[\cG;\bar\cA]\;\cW(\x,\y,\z;\bx,\byp,\byp)
\nonumber\\
&\qquad+\cH[\cA;\bar\cA]\;\cW(\x,\yp,\yp;\bx,\byp,\byp)\Big]\,.
\label{eq:xs-qbar}
\end{align}
The quark mass depends on the flavour, so the sum over flavours acts on the whole integrand. The
function $\cH$ is the sum over the helicities of the quark and the antiquark and over the polarisations
of the photon and of the gluon of the product of the effective vertices,
\begin{equation}
\cH[F;\bar F]\equiv\sum_{\lambda_1,\lambda_2}\sum_{l,r}\cM^{lr}_{\lambda_1\lambda_2}[F](\x,\y,\z)\;
\cM^{lr\,*}_{\lambda_1\lambda_2}[\bar F](\bx,\by,\bz)\,,
\label{eq:Hdef}
\end{equation}
where $F$ stands for either $\cG$ or $\cA$, with $\cM[\cG]\equiv\cM_1$ and $\cM[\cA]\equiv\cM_2$ of
eq.~\eqref{eq:M}. To treat the two cases on the same footing we associate to each set a fifth function
$F_E$, which multiplies the instantaneous term: $F_E=\cG_E$ for $F=\cG$ and $F_E=0$ for $F=\cA$. We
find $\cH[F;\bar F]=\cH_{\rm reg\times reg}+\cH_{\rm reg\times inst}+\cH_{\rm inst\times inst}$, where the
three terms are the following functionals of $F$ and $\bar F$:
\begin{subequations}
\label{eq:H-direct}
\begin{align}
\cH_{\rm reg\times reg}={}&\Big\{8\big[\zeta^2+(1-\zeta)^2\big]\big[\eta^2+(1-\zeta)^2\big]\,
(\R\cdot\bR)(\Y\cdot\bY)
\nonumber\\
&\quad-8\,\xi\,(2\zeta-1)\big(2(1-\zeta)-\xi\big)\,(\R\times\bR)(\Y\times\bY)\Big\}\,F^{(11)}\bar F^{(11)}
\nonumber\\
&+8\,\mf^2\big[\eta^2+(1-\zeta)^2\big]\,(\Y\cdot\bY)\,F^{(01)}\bar F^{(01)}
+8\,\mf^2\,\xi^2\of^2\,F^{(00)}\bar F^{(00)}
\nonumber\\
&+8\,\xi^2\of^2\big[\zeta^2+(1-\zeta)^2\big]\,(\R\cdot\bR)\,F^{(10)}\bar F^{(10)}
\nonumber\\
&-8\,\zeta\,\eta\,\mf\,\xi\,\of\Big[(\R\cdot\Y)\,F^{(11)}\bar F^{(00)}+(\bR\cdot\bY)\,F^{(00)}\bar F^{(11)}
\nonumber\\
&\hspace{2.6cm}+(\R\cdot\bY)\,F^{(10)}\bar F^{(01)}+(\Y\cdot\bR)\,F^{(01)}\bar F^{(10)}\Big]\,,
\label{eq:Hrr}\\[1mm]
\cH_{\rm reg\times inst}={}&-\frac{8\,\zeta\,\xi\,\eta^2}{1-\zeta}\,
\Big[(\R\cdot\Y)\,F^{(11)}\bar F_E+(\bR\cdot\bY)\,F_E\,\bar F^{(11)}\Big]
\nonumber\\
&+\frac{8\,\xi^2\,\eta\,\mf\,\of}{1-\zeta}\,\Big[F^{(00)}\bar F_E+F_E\,\bar F^{(00)}\Big]\,,
\label{eq:Hri}\\[1mm]
\cH_{\rm inst\times inst}={}&8\left[\frac{\xi\,\eta}{1-\zeta}\right]^2F_E\,\bar F_E\,,
\label{eq:Hii}
\end{align}
\end{subequations}
where $\bm{A}\times\bm{B}\equiv\ee^{ij}A^iB^j$. The first two lines of eq.~\eqref{eq:Hrr}, together
with the first line of eq.~\eqref{eq:Hri} and with eq.~\eqref{eq:Hii}, are the massless structures. The
mass generates (i) three ``diagonal'' terms, in which the same helicity path is followed in the DA and
in the CCA, with weights given by the squares of the photon splitting factor
$\zeta^2+(1-\zeta)^2$ or of the gluon splitting factor $\eta^2+(1-\zeta)^2$ and by the squares of the
flip vertices $\mf$ and $\xi\of$; (ii) four terms in which two different paths interfere, namely the
no-flip path with the double-flip path and the two single-flip paths with each other, all with the
common weight $-8\zeta\eta\,\mf\,\xi\of$; and (iii) the interference of the double-flip path with the
instantaneous term. Paths with a different total number of flips do not interfere, since they lead to
different final helicities.

\subsection{Direct contribution of the quark}
\label{sec:xs-q}

Similarly, from eq.~\eqref{eq:state-q},
\begin{align}
\frac{\dd\sigma^{(qq)}}{\dd^3k_1\,\dd^3k_2\,\dd^3k_3}=&\;
\frac{\aem\,\as\,\CF\Nc}{2(2\pi)^{10}}\;\frac{\eta^2}{\xi\,(q^+)^2}\;
\delta\big(q^+-k_1^+-k_2^+-k_3^+\big)\sum_fe_f^2
\nonumber\\
&\times\int_{\x,\y,\z,\bx,\by,\bz}e^{-i\kk_1\cdot(\x-\bx)-i\kk_2\cdot(\y-\by)-i\kk_3\cdot(\z-\bz)}\;
\Big[\tH[\tG;\bar\tG]\;\cW(\x,\y,\z;\bx,\by,\bz)
\nonumber\\
&\qquad-\tH[\tA;\bar\tG]\;\cW(\xp,\y,\xp;\bx,\by,\bz)
-\tH[\tG;\bar\tA]\;\cW(\x,\y,\z;\bxp,\by,\bxp)
\nonumber\\
&\qquad+\tH[\tA;\bar\tA]\;\cW(\xp,\y,\xp;\bxp,\by,\bxp)\Big]\,,
\label{eq:xs-q}
\end{align}
with $\tH[\tF;\bar\tF]\equiv\sum\tM^{lr}_{\lambda_1\lambda_2}[\tF]\,\tM^{lr\,*}_{\lambda_1\lambda_2}[\bar\tF]$,
$\tF_E=\cG_E$ for $\tF=\tG$ and $\tF_E=0$ for $\tF=\tA$, and with
$\tH[\tF;\bar\tF]=\tH_{\rm reg\times reg}+\tH_{\rm reg\times inst}+\tH_{\rm inst\times inst}$,
\begin{subequations}
\label{eq:Ht-direct}
\begin{align}
\tH_{\rm reg\times reg}={}&\Big\{8\big[\zeta^2+(1-\eta)^2\big]\big[\eta^2+(1-\eta)^2\big]\,
(\Rt\cdot\bRt)(\X\cdot\bX)
\nonumber\\
&\quad-8\,\xi\,(2\eta-1)\big(2(1-\eta)-\xi\big)\,(\Rt\times\bRt)(\X\times\bX)\Big\}\,\tF^{(11)}\bar\tF^{(11)}
\nonumber\\
&+8\,\mf^2\big[\zeta^2+(1-\eta)^2\big]\,(\X\cdot\bX)\,\tF^{(01)}\bar\tF^{(01)}
+8\,\mf^2\,\xi^2\oft^{\,2}\,\tF^{(00)}\bar\tF^{(00)}
\nonumber\\
&+8\,\xi^2\oft^{\,2}\big[\eta^2+(1-\eta)^2\big]\,(\Rt\cdot\bRt)\,\tF^{(10)}\bar\tF^{(10)}
\nonumber\\
&+8\,\zeta\,\eta\,\mf\,\xi\,\oft\Big[(\Rt\cdot\X)\,\tF^{(11)}\bar\tF^{(00)}+(\bRt\cdot\bX)\,\tF^{(00)}\bar\tF^{(11)}
\nonumber\\
&\hspace{2.6cm}+(\Rt\cdot\bX)\,\tF^{(10)}\bar\tF^{(01)}+(\X\cdot\bRt)\,\tF^{(01)}\bar\tF^{(10)}\Big]\,,
\label{eq:Htrr}\\[1mm]
\tH_{\rm reg\times inst}={}&\frac{8\,\zeta^2\,\xi\,\eta}{1-\eta}\,
\Big[(\Rt\cdot\X)\,\tF^{(11)}\bar\tF_E+(\bRt\cdot\bX)\,\tF_E\,\bar\tF^{(11)}\Big]
\nonumber\\
&+\frac{8\,\xi^2\,\zeta\,\mf\,\oft}{1-\eta}\,\Big[\tF^{(00)}\bar\tF_E+\tF_E\,\bar\tF^{(00)}\Big]\,,
\label{eq:Htri}\\[1mm]
\tH_{\rm inst\times inst}={}&8\left[\frac{\xi\,\zeta}{1-\eta}\right]^2\tF_E\,\bar\tF_E\,.
\label{eq:Htii}
\end{align}
\end{subequations}
These expressions follow from eqs.~\eqref{eq:H-direct} by the substitutions
$\zeta\leftrightarrow\eta$, $\R\to-\Rt$, $\Y\to\X$ and $\of\to\oft$; the terms which are odd in $\R$
change sign.

\subsection{Interference contribution}
\label{sec:xs-int}

The interference between an emission by the quark in the DA and an emission by the antiquark in the
CCA, plus its complex conjugate, gives
\begin{align}
\frac{\dd\sigma^{(q\bar q)}}{\dd^3k_1\,\dd^3k_2\,\dd^3k_3}=&
-\frac{\aem\,\as\,\CF\Nc}{2(2\pi)^{10}}\;\frac{\zeta\,\eta}{\xi\,(q^+)^2}\;
\delta\big(q^+-k_1^+-k_2^+-k_3^+\big)\sum_fe_f^2
\nonumber\\
&\times2\,\re\int_{\x,\y,\z,\bx,\by,\bz}e^{-i\kk_1\cdot(\x-\bx)-i\kk_2\cdot(\y-\by)-i\kk_3\cdot(\z-\bz)}
\nonumber\\
&\times\Big[\cH^\times[\tG;\bar\cG]\;\cW(\x,\y,\z;\bx,\by,\bz)
\nonumber\\
&\qquad-\cH^\times[\tA;\bar\cG]\;\cW(\xp,\y,\xp;\bx,\by,\bz)
-\cH^\times[\tG;\bar\cA]\;\cW(\x,\y,\z;\bx,\byp,\byp)
\nonumber\\
&\qquad+\cH^\times[\tA;\bar\cA]\;\cW(\xp,\y,\xp;\bx,\byp,\byp)\Big]\,,
\label{eq:xs-int}
\end{align}
where the overall sign comes from the opposite signs of eqs.~\eqref{eq:state-qbar} and
\eqref{eq:state-q}, and
\begin{equation}
\cH^\times[\tF;\bar F]\equiv\sum_{\lambda_1,\lambda_2}\sum_{l,r}\tM^{lr}_{\lambda_1\lambda_2}[\tF](\x,\y,\z)\;
\cM^{lr\,*}_{\lambda_1\lambda_2}[\bar F](\bx,\by,\bz)\,.
\end{equation}
We find $\cH^\times[\tF;\bar F]=\cH^\times_{\rm reg\times reg}+\cH^\times_{\rm reg\times inst}
+\cH^\times_{\rm inst\times inst}$, with
\begin{subequations}
\label{eq:Hx}
\begin{align}
\cH^\times_{\rm reg\times reg}={}&\Big\{8\big[1-\xi-2\zeta\eta\big]\big[\xi(1-\xi)+2\zeta\eta\big]\,
(\Rt\cdot\bR)(\X\cdot\bY)
\nonumber\\
&\quad-8\,\xi\,\big[1-2\zeta-\xi\big]^2\,(\Rt\times\bR)(\X\times\bY)\Big\}\,\tF^{(11)}\bar F^{(11)}
\nonumber\\
&+8\,\mf^2\,(\xi+2\zeta\eta)\,(\X\cdot\bY)\,\tF^{(01)}\bar F^{(01)}
-8\,(\xi+2\zeta\eta)\,\xi^2\,\oft\,\of\,(\Rt\cdot\bR)\,\tF^{(10)}\bar F^{(10)}
\nonumber\\
&-8\,(\zeta+\xi)^2\,\mf\,\xi\,\of\,(\Rt\cdot\X)\,\tF^{(11)}\bar F^{(00)}
+8\,(1-\zeta)^2\,\mf\,\xi\,\oft\,(\bR\cdot\bY)\,\tF^{(00)}\bar F^{(11)}
\nonumber\\
&+8\,\eta^2\,\mf\,\xi\,\oft\,(\Rt\cdot\bY)\,\tF^{(10)}\bar F^{(01)}
-8\,\zeta^2\,\mf\,\xi\,\of\,(\X\cdot\bR)\,\tF^{(01)}\bar F^{(10)}\,,
\label{eq:Hxrr}\\[1mm]
\cH^\times_{\rm reg\times inst}={}&-\frac{8\,\xi\,\eta\,(\zeta+\xi)^2}{1-\zeta}\,
(\Rt\cdot\X)\,\tF^{(11)}\bar F_E
\nonumber\\
&+\frac{8\,\xi\,\zeta\,(1-\zeta)^2}{\zeta+\xi}\,(\bR\cdot\bY)\,\tF_E\,\bar F^{(11)}\,,
\label{eq:Hxri}\\[1mm]
\cH^\times_{\rm inst\times inst}={}&0\,.
\label{eq:Hxii}
\end{align}
\end{subequations}
Two features of the interference term are worth noticing. First, the two instantaneous terms do not
interfere with each other: the one of the antiquark is proportional to
$\delta^{lr}+2i\ee^{lr}\lambda_1$ and the one of the quark to $\delta^{lr}-2i\ee^{lr}\lambda_1$, and
$\sum_{l,r}(\delta^{lr}-2i\ee^{lr}\lambda_1)(\delta^{lr}+2i\ee^{lr}\lambda_1)^*=2-8\lambda_1^2=0$. Second,
the double-flip path of one emitter interferes neither with the instantaneous term nor with the
double-flip path of the other emitter, because the two flip vertices appear in the opposite order
along the fermion line; the double-flip paths contribute to the interference only through their
overlap with the no-flip path.

Equations \eqref{eq:xs-qbar}, \eqref{eq:xs-q} and \eqref{eq:xs-int}, together with the helicity sums
\eqref{eq:H-direct}, \eqref{eq:Ht-direct} and \eqref{eq:Hx}, the radial functions
\eqref{eq:Gab}, \eqref{eq:Aexplicit}, \eqref{eq:GE}, \eqref{eq:Gtab} and \eqref{eq:Atexplicit}, and the
colour structures of section~\ref{sec:colour}, are our complete result for the leading-order trijet
cross-section with massive quarks, for a transversely polarised virtual photon.
